\SGARefTarget{sga3:XXIII:2.5}
\SGAStatement{Corollary 2.5}{Let \(S\) be a scheme; let \(G\) and
\(G'\) be two pinned \(S\)-groups of semisimple rank \(2\); let
\(q>0\) be an integer such that \(x\mapsto x^q\) is an endomorphism
of \(\mathbf G_{a,S}\); and let
\[
  h:R(G')\longrightarrow R(G)
\]
be a \(q\)-morphism of pinned root data. For each
\(\alpha\in R^+\), choose
\[
  X_\alpha\in\Gamma(S,\mathfrak g^\alpha)^\times,
  \qquad
  X'_{d(\alpha)}
    \in\Gamma(S,\mathfrak g'^{\,d(\alpha)})^\times,
\]
extending the canonical choices for \(\alpha\in\Delta\). Suppose:
\begin{enumerate}
  \item[\SGAInlineRefTarget{sga3:XXIII:2.5:item:i}(i)] if \(\Delta=\{\alpha,\beta\}\), then
    \[
      D_S(h)(t_{\alpha\beta})=t'_{d(\alpha)d(\beta)};
    \]
  \item[\SGAInlineRefTarget{sga3:XXIII:2.5:item:ii}(ii)] for every \(\alpha\in\Delta\) and every
    \(\beta\in R^+\), \(\beta\ne\alpha\), if
    \(z\in\mathbf G_m(S)\) is defined by
    \[
      \operatorname{Ad}(w_\alpha)X_\beta
        =zX_{s_\alpha(\beta)},
    \]
    then
    \[
      \operatorname{Ad}(w'_{d(\alpha)})X'_{d(\beta)}
        =z^{q(\beta)}X'_{d(s_\alpha(\beta))};
    \]
  \item[(iii)] there is a morphism of groups \(f:U\to U'\) such
    that, for every \(\alpha\in R^+\), every \(S'\to S\), and every
    \(x\in\mathbf G_a(S')\),
    \[
      f(\exp(xX_\alpha))
        =\exp\!\bigl(x^{q(\alpha)}X'_{d(\alpha)}\bigr).
    \]
\end{enumerate}
Then there is a morphism of pinned groups \(g:G\to G'\) such that
\(R(g)=h\).}

Indeed, define \(f_\alpha:U_\alpha\to G'\) by
\[
  f_\alpha(\exp(xX_\alpha))
    =\exp\!\bigl(x^{q(\alpha)}X'_{d(\alpha)}\bigr),
\]
and put
\[
  f_T=D_S(h),\qquad h_\alpha=w'_{d(\alpha)}.
\]
The conditions of \SGARefLink{sga3:XXIII:2.3}{Theorem 2.3} hold: note that
\[
  q(s_\alpha(\beta))=q(\beta)
\]
by \SGARefLink{sga3:XXI:6.8.4}{Exposé XXI, 6.8.4}, and that
\[
  D_S(h)(t_\alpha)=t'_{d(\alpha)}
\]
always holds. The assertion follows at once.

\SGARefTarget{sga3:XXIII:2.6}
\SGAStatement{Remark 2.6}{Condition (v) of \SGARefLink{sga3:XXIII:2.3}{Theorem 2.3} may be made
more explicit as follows. One must first verify:
\begin{enumerate}
  \item[\SGAInlineRefTarget{sga3:XXIII:2.6:item:a}(a)] for every word in \(w_\alpha,w_\beta\) whose
    corresponding word carries \(\alpha\) or \(\beta\) to
    \(\alpha\) or \(\beta\), the corresponding relation of the type
    occurring in \SGARefLink{sga3:XXIII:2.3.3}{2.3.3} holds.
\end{enumerate}
The proof of \SGARefLink{sga3:XXIII:2.3.3}{2.3.3} shows that it suffices to check the words that are
minimal, in the sense that no nontrivial initial subword satisfies
the imposed conditions.

If \SGARefLink{sga3:XXIII:2.6:item:a}{(a)} holds, define \(f_\gamma:U_\gamma\to H\) for every
\(\gamma\in R^+_{\alpha\beta}\) as in \SGARefLink{sga3:XXIII:2.3.5}{2.3.5}. One must then require:
\begin{enumerate}
  \item[\SGAInlineRefTarget{sga3:XXIII:2.6:item:b}(b)] the map defined by the \(f_\gamma\),
    \[
      U_{\alpha\beta}
        =\prod_{\gamma\in R^+_{\alpha\beta}}U_\gamma
        \longrightarrow H,
    \]
    is a morphism of groups.
\end{enumerate}
By \SGARefLink{sga3:xxii:lem:5.5.8}{Exposé XXII, 5.5.8}, condition \SGARefLink{sga3:XXIII:2.6:item:b}{(b)} follows from:
\begin{enumerate}
  \item[(b\('\))] the preceding map respects the commutation
    relations between \(U_\gamma\) and \(U_\delta\), for
    \(\gamma,\delta\in R^+_{\alpha\beta}\), \(\gamma>\delta\);
    namely, the relations involving
    \(C_{i,j,\gamma,\delta}\) in \SGARefLink{sga3:xxii:cor:5.5.2}{Exposé XXII, 5.5.2}.
\end{enumerate}}

Conversely, the conjunction of \SGARefLink{sga3:XXIII:2.6:item:a}{(a)} and (b\('\)) is equivalent to
condition (v).

These conditions may even be reduced to conditions involving only
\[
  h_\alpha,\quad h_\beta,\quad
  f_\alpha(u_\alpha),\quad f_\beta(u_\beta)
\]
in \(H\). For example, if
\(s_\alpha s_\beta s_\alpha(\beta)=\alpha\), a condition of type
(a) reads
\begin{equation}
  \operatorname{int}(h_\alpha h_\beta h_\alpha)f_\beta(x)
    =f_\alpha\!\left(
       \operatorname{int}(w_\alpha w_\beta w_\alpha)x
     \right).
  \tag{1}
  \label{sga3:XXIII:2.6:eq1}
\end{equation}
Here \(x\in U_\beta(S')\), \(S'\to S\). In particular, for
\(x=u_\beta\),
\[
  \operatorname{int}(w_\alpha w_\beta w_\alpha)u_\beta
    =u_\alpha^z
\]
for a certain \(z\in\mathbf G_m(S)\), and
\begin{equation}
  \operatorname{int}(h_\alpha h_\beta h_\alpha)f_\beta(u_\beta)
    =f_\alpha(u_\alpha)^z.
  \tag{1\('\)}
  \label{sga3:XXIII:2.6:eq1prime}
\end{equation}

Conversely, taking conditions \SGARefLink{sga3:XXIII:2.5:item:i}{(i)}--(iv) of \SGARefLink{sga3:XXIII:2.3}{2.3} into account,
\eqref{sga3:XXIII:2.6:eq1prime} implies
\eqref{sga3:XXIII:2.6:eq1}. If \(t\in T(S')\), conjugating
\eqref{sga3:XXIII:2.6:eq1prime} by \(t\), and using \SGARefLink{sga3:XXIII:2.3:item:i}{(i)} and \SGARefLink{sga3:XXIII:2.3:item:ii}{(ii)},
gives \eqref{sga3:XXIII:2.6:eq1} for
\[
  x=\operatorname{int}(t)u_\beta=u_\beta^{\beta(t)}.
\]
The morphism \(\beta:T\to\mathbf G_{m,S}\) is faithfully flat, so
\eqref{sga3:XXIII:2.6:eq1} holds for
\(x\in U_\beta(S')^\times\). It is additive in \(x\), and every
section of \(U_\beta\) is locally the sum of two invertible sections.
Thus
\[
  \text{(1')}+\text{(i)}+\text{(ii)}
    \quad\Longrightarrow\quad\text{(1)}.
\]

The same argument applies to conditions of type \SGARefLink{sga3:XXIII:2.6:item:b}{(b)}. One uses the
fact that, for two distinct positive roots \(\gamma,\gamma'\), hence
linearly independent over \(\mathbf Q\), the morphism
\[
  T\longrightarrow\mathbf G_{m,S}^2
\]
with components \(\gamma,\gamma'\) is faithfully flat. We leave the
details of this transposition to the reader.

\SGANextSectionTarget{sga3:xxiii:section:3}\SGASection{3}{Groups of Semisimple Rank \(2\)}
\SGANextSectionTarget{sga3:xxiii:subsection:3-1}\SGASubsection{3.1}{Generalities}

\SGARefTarget{sga3:XXIII:3.1.1}
\SGAStatement{Lemma 3.1.1}{Let \(S\) be a scheme, let \(G\) be a
pinned \(S\)-group, and let \(\alpha,\beta\) be two roots of \(G\)
such that \(\alpha+\beta\ne0\).
\begin{enumerate}
  \item[\SGAInlineRefTarget{sga3:xxiii:subsection:3-1:item:i}\SGAInlineRefTarget{sga3:xxiii:lemma:3-1-1:item:i}(i)] If \(\alpha+\beta\notin R\), then
    \[
      \exp(X_\alpha)\exp(X_\beta)
        =\exp(X_\beta)\exp(X_\alpha)
    \]
    for every \(S'\to S\),
    \(X_\alpha\in W(\mathfrak g^\alpha)(S')\), and
    \(X_\beta\in W(\mathfrak g^\beta)(S')\).
  \item[\SGAInlineRefTarget{sga3:xxiii:lemma:3-1-1:item:ii}\SGAInlineRefTarget{sga3:xxiii:subsection:3-1:item:ii}(ii)] If neither \(\alpha+\beta\) nor \(\alpha-\beta\) is a
    root, then
    \[
      w_\alpha(z_\alpha)\exp(X_\beta)w_\alpha(z_\alpha)^{-1}
        =\exp(X_\beta)
    \]
    for every \(X_\beta\in W(\mathfrak g^\beta)(S')\) and
    \(z_\alpha\in W(\mathfrak g^\alpha)^\times(S')\), and
    \[
      w_\alpha(z_\alpha)w_\beta(z_\beta)
        =w_\beta(z_\beta)w_\alpha(z_\alpha)
    \]
    for every
    \(z_\alpha\in W(\mathfrak g^\alpha)^\times(S')\) and
    \(z_\beta\in W(\mathfrak g^\beta)^\times(S')\).
  \item[\SGAInlineRefTarget{sga3:XXIII:3.1.1:item:iii}\SGAInlineRefTarget{sga3:xxiii:lemma:3-1-1:item:iii}(iii)] Let
    \(X_\alpha\in W(\mathfrak g^\alpha)^\times(S')\),
    \(X_\beta\in W(\mathfrak g^\beta)^\times(S')\), and
    \(w\in\operatorname{Norm}_G(T)(S')\) satisfy
    \(w(\alpha)=\beta\). Define \(z\in\mathbf G_m(S')\) by
    \[
      \operatorname{Ad}(w)X_\alpha=zX_\beta.
    \]
    Then
    \[
    \begin{aligned}
      \operatorname{int}(w)\exp(xX_\alpha)
        &=\exp(xzX_\beta),\\
      \operatorname{int}(w)\exp(yX_\alpha^{-1})
        &=\exp(yz^{-1}X_\beta^{-1}),\\
      \operatorname{int}(w)w_\alpha(X_\alpha)
        &=\beta^*(z)w_\beta(X_\beta).
    \end{aligned}
    \]
    In particular, if \(z=\pm1\), then
    \[
      \operatorname{int}(w)w_\alpha(X_\alpha)
        =w_\beta(X_\beta)^z.
    \]
  \item[\SGAInlineRefTarget{sga3:XXIII:3.1.1:item:iv}(iv)] If
    \(t_\alpha=\alpha^*(-1)\) and \(t_\beta=\beta^*(-1)\), then
    \[
      s_\alpha(t_\beta)
        =t_\beta t_\alpha^{(\beta^*,\alpha)},
      \qquad
      \beta(t_\alpha)=(-1)^{(\alpha^*,\beta)}.
    \]
\end{enumerate}}

\SGAProof
Part \SGARefLink{sga3:xxiii:subsection:3-1:item:i}{(i)} follows immediately from \SGARefLink{sga3:xxii:cor:5.5.2}{Exposé XXII, 5.5.2}. Part \SGARefLink{sga3:xxiii:subsection:3-1:item:ii}{(ii)}
follows from \SGARefLink{sga3:xx:3.1}{Exposé XX, 3.1} and from \SGARefLink{sga3:xxiii:subsection:3-1:item:i}{(i)} applied to
\[
  (\beta,\alpha),\quad(\beta,-\alpha),\quad
  (-\beta,\alpha),\quad(-\beta,-\alpha).
\]
Part \SGARefLink{sga3:XXIII:3.1.1:item:iii}{(iii)} is immediate from the definitions; for its last assertion,
observe that
\[
  \beta^*(-1)=w_\beta(X_\beta)^{-2}.
\]
Finally, \SGARefLink{sga3:XXIII:3.1.1:item:iv}{(iv)} is tautological.

\SGARefTarget{sga3:XXIII:3.1.2}
\SGAStatement{Proposition 3.1.2 (Groups of type
\(A_1\times A_1\))}{Let \(S\) be a scheme and \(G\) a pinned
\(S\)-group of type \(A_1\times A_1\). Write
\(\Delta=R^+=\{\alpha,\beta\}\).
\begin{enumerate}
  \item[\SGAInlineRefTarget{sga3:XXIII:3.1.2:item:i}(i)]
    \[
      t_{\alpha\beta}=(w_\alpha w_\beta)^2
        =t_\alpha t_\beta
        =(w_\beta w_\alpha)^2=t_{\beta\alpha}.
    \]
  \item[\SGAInlineRefTarget{sga3:XXIII:3.1.2:item:ii}(ii)]\SGAInlineRefTarget{sga3:xxiii:proposition:3-1-2:item:ii}
    \[
      \operatorname{Ad}(w_\alpha)X_\beta=X_\beta,
      \qquad
      \operatorname{Ad}(w_\beta)X_\alpha=X_\alpha.
    \]
  \item[\SGAInlineRefTarget{sga3:XXIII:3.1.2:item:iii}(iii)] The groups \(U_\alpha\) and \(U_\beta\) commute; in
    other words, \(U\) is commutative.
\end{enumerate}}

\SGAProof
\SGARefLink{sga3:xxiii:lemma:3-1-1:item:ii}{Lemma 3.1.1(ii)} gives \(w_\alpha w_\beta=w_\beta w_\alpha\), whence
\[
  (w_\alpha w_\beta)^2=w_\alpha^2w_\beta^2=t_\alpha t_\beta,
\]
which proves \SGARefLink{sga3:XXIII:3.1.2:item:i}{(i)}. The same lemma gives \SGARefLink{sga3:XXIII:3.1.2:item:ii}{(ii)}, and \SGARefLink{sga3:XXIII:3.1.2:item:iii}{(iii)} is
\SGARefLink{sga3:xxiii:lemma:3-1-1:item:i}{Lemma 3.1.1(i)}.

\SGARefTarget{sga3:XXIII:3.1.3}
\SGAParagraph{3.1.3}
Condition (v) of \SGARefLink{sga3:XXIII:2.3}{Theorem 2.3} takes the following explicit form. Put
\[
  v_\alpha=f_\alpha(u_\alpha)
  \qquad(\alpha\in\Delta).
\]
Then the two groups of conditions are
\[
\text{(A)}\quad
\begin{cases}
  h_\alpha v_\beta h_\alpha^{-1}=v_\beta,\\
  h_\beta v_\alpha h_\beta^{-1}=v_\alpha,
\end{cases}
\qquad
\text{(B)}\quad
v_\alpha v_\beta=v_\beta v_\alpha.
\]
